% Hodina 48, 1. ročník SPŠ. Rozsah: Rozpis_vyuky_MAT_1_rocnik.tex, řádek 48.

% Plán: 1_rocnik/tematicke_plany/SPS.pdf; společná projektová šablona.

% Učebnice: ucebnice/01_Zakladni_poznatky/Algebraicke_vyrazy/Algebraické výrazy.tex

% Původ úloh: P2: učebnice, „Rozklad mnohočlenů“, příklad se třemi proměnnými. Ostatní vlastní.

% Metodický zdroj: externi_zdroje/realisticky/Mnohocleny/06_Rozklad_mnohoclenu_na_soucin_I_vytykani_.pdf

% URL: http://www.realisticky.cz/ucebnice/01%20Matematika%20S%C5%A0/01%20Z%C3%A1kladn%C3%AD%20poznatky/08%20Mnoho%C4%8Dleny/06%20Rozklad%20mnoho%C4%8Dlen%C5%AF%20na%20sou%C4%8Din%20I%20(vyt%C3%BDk%C3%A1n%C3%AD).pdf

\documentclass[11pt,a4paper]{article}

% Spolecny styl pro ucitelske pripravy MAT (Work/MAT)
% Verze 1.1, 2026-10-23
\NeedsTeXFormat{LaTeX2e}
\ProvidesPackage{priprava}[2026/10/23 v1.1 Sablona ucitelskych priprav MAT]

% Zakladni typografie a cestina
\RequirePackage[T1]{fontenc}
\RequirePackage[utf8]{inputenc}
\RequirePackage[czech]{babel}
\RequirePackage{lmodern}
\RequirePackage{microtype}
\RequirePackage{amsmath,amssymb}
\RequirePackage{xcolor}
\RequirePackage{enumitem}
\RequirePackage[most]{tcolorbox}
\RequirePackage{tabularx,booktabs,array}
\RequirePackage{fancyhdr}
\RequirePackage{needspace}
\RequirePackage{geometry}
\RequirePackage[hidelinks]{hyperref}

\geometry{a4paper,top=17mm,bottom=18mm,left=23mm,right=23mm}
\definecolor{notegray}{gray}{0.39}
\setlength{\parindent}{0pt}
\setlength{\parskip}{2.5pt plus 1pt minus 0.5pt}
\setlist[itemize]{topsep=2pt,itemsep=1pt,parsep=0pt,leftmargin=1.45em}
\setlist[enumerate]{topsep=2pt,itemsep=2pt,parsep=0pt,leftmargin=1.65em}
\setcounter{secnumdepth}{0}
\renewcommand{\arraystretch}{1.12}

% Zahlavi: vlevo rocnik | globalni tema, vpravo cislo: nazev hodiny
% Zapati zustava zcela prazdne.
\pagestyle{fancy}
\fancyhf{}
\renewcommand{\headrulewidth}{0.35pt}
\renewcommand{\footrulewidth}{0pt}
\setlength{\headheight}{14pt}
\newcommand{\nastavPripravu}[4]{%
  \fancyhead[L]{\small #1\enspace\textbar\enspace #2}%
  \fancyhead[R]{\small #3: #4}%
}

% Minimalisticke cernobile boxy. Kazda definice zvlast.
\tcbset{commonlesson/.style={enhanced,breakable=false,colback=white,colframe=black,
  colbacktitle=white,coltitle=black,boxrule=0.45pt,
  arc=0.8mm,outer arc=0.8mm,left=3mm,right=3mm,top=1.6mm,bottom=1.6mm,
  toptitle=1mm,bottomtitle=0.7mm,fonttitle=\bfseries,
  attach boxed title to top left={xshift=2.5mm,yshift=-1.15mm},
  boxed title style={colback=white,colframe=white,boxrule=0pt,left=1mm,right=1mm}}}
\newtcolorbox{zadaniBox}[1][]{commonlesson,title={Zadání},#1}
\newtcolorbox{prikladBox}[1][]{commonlesson,title={Příklad},#1}
\newtcolorbox{definiceBox}[1][]{commonlesson,title={Definice},#1}

% Spolecny box pro zadani a jeho reseni, bez mezery mezi castmi.
% Pouziti:
%   \begin{prikladReseniBox}
%     Zadani ulohy.
%     \tcblower
%     \begin{align*} ... \end{align*}
%   \end{prikladReseniBox}
% Spodni cast ma jemny sedy podklad a delici linku; lze v ni sazet
% vice radku matematiky. Obe casti maji jen jeden spolecny ramecek.
\tcbset{solutionlesson/.style={
  commonlesson,
  enhanced jigsaw,
  breakable,
  colbacklower=black!3,
  segmentation style={solid,draw=black!25,line width=0.35pt},
  fontlower=\small\color{notegray},
  before lower={\noindent\textbf{Řešení:}\enspace},
  leftlower=3mm,rightlower=3mm
}}
\newtcolorbox{prikladReseniBox}[1][]{solutionlesson,title={Příklad},#1}
\newtcolorbox{zadaniReseniBox}[1][]{solutionlesson,title={Zadání},#1}

% Pomucky pro autora pripravy: sedive poznamky; \reseni pro starsi pripravy
\newcommand{\poznamka}[1]{\par\begingroup\small\color{notegray}\itshape
  \textnormal{\textbf{Poznámka pro učitele.}}\enspace #1\par\endgroup}
\newcommand{\reseni}[1]{\par\begingroup\small\color{notegray}
  \textbf{Očekávané řešení:}\enspace #1\par\endgroup}
\newcommand{\blok}[1]{\needspace{5\baselineskip}\subsection*{#1}}
\newcommand{\vyrok}[1]{\uv{#1}}


\nastavPripravu{1. ročník}{Algebraické výrazy}{48}{Vytýkání společného činitele}

\begin{document}

\blok{Výklad a souvislosti}

Dosud jsme součin roznásobovali. Nyní budeme hledat společné činitele jednotlivých členů a psát výraz opačným směrem. Správnost každého rozkladu lze ověřit zpětným roznásobením.

\begin{definiceBox}[title={Definice},breakable=false]
\textbf{Rozklad mnohočlenu na součin} je jeho nahrazení součinem mnohočlenů, který má pro všechna dosazení stejnou hodnotu.
\end{definiceBox}


\begin{definiceBox}[title={Definice},breakable=false]
\textbf{Vytýkání} je použití distributivního zákona v opačném směru: $MA+MB=M(A+B)$. Společný činitel může být jednočlen i celá závorka.
\end{definiceBox}


U číselných koeficientů hledáme společný dělitel, u proměnných nejnižší společný exponent. Po vytknutí vydělíme společným činitelem každý člen. Když společný činitel nemají všechny členy, může pomoci seskupení.

\blok{Řešené příklady}

\begin{prikladReseniBox}[title={Příklad 1},breakable=false]
Rozlož $6x^3-9x^2$.
\tcblower
Oba členy obsahují $3x^2$. Dostáváme $3x^2(2x-3)$. Roznásobení vrací původní dva členy; jako mnohočlenná identita platí i pro $x=0$.
\end{prikladReseniBox}


\begin{prikladReseniBox}[title={Příklad 2},breakable=false]
Rozlož $7a^3bc-14a^3b^2c^2-21ab^4$.
\tcblower
Vytkneme $7ab$ a zůstane $7ab(a^2c-2a^2bc^2-3b^3)$. Každý člen v závorce musí po násobení $7ab$ dát právě jeden původní člen.
\end{prikladReseniBox}


\blok{Samostatné procvičení}

\begin{zadaniReseniBox}[title={Procvičení},breakable=false]
\begin{enumerate}[label=\alph*)]
\item Rozlož $12a^2b-18ab^2$.
\item Rozlož $-3x^2+6x$.
\item Rozlož $x^3+2x^2+3x+6$ seskupením.
\end{enumerate}
\tcblower
\begin{enumerate}[label=\alph*)]
\item $6ab(2a-3b)$. Žádné zákazy $a=0$ ani $b=0$ nevznikají, protože výsledkem je identita bez dělení.
\item $-3x(x-2)$; vytažení záporného koeficientu změní znaménka v závorce.
\item $x^2(x+2)+3(x+2)=(x+2)(x^2+3)$.
\end{enumerate}
\end{zadaniReseniBox}


\poznamka{Nedoplňuj automaticky podmínku nenulovosti vytaženého činitele: na rozdíl od krácení zlomku jde o mnohočlennou identitu. Při pomalejším tempu řeš seskupování společně. Vzorce pro rozdíl čtverců zde zatím nepoužívej.}

\end{document}
